\answer{OP22}
While the pattern is still valid and the current value is $x$, let $f(x)$ be the expected number of additional draws when the next comparison must be an increase; let $g(x)$ be the corresponding expectation when a decrease is required. One draw is always made, and continuation occurs only if that comparison succeeds. Thus
\[
f(x)=1+\int_x^1g(y)\dd y,\qquad
 g(x)=1+\int_0^xf(y)\dd y.
\]
The expectations are finite. If the first $n$ values obey the pattern, every disjoint pair $(U_1,U_2),(U_3,U_4),\ldots$ must increase, so the survival probability is at most $2^{-\lfloor n/2\rfloor}$. The same geometric bound controls the remaining draws from any state and justifies the integral equations.

Differentiating gives $f'=-g$ and $g'=f$, hence $f''+f=0$. The boundary conditions $f(1)=1$ and $g(0)=1$, or $f'(0)=-1$, yield
\[
f(x)=A\cos x-\sin x,\qquad
A=\frac{1+\sin1}{\cos1}=\sec1+\tan1.
\]
After the first draw, the next comparison must be an increase. Therefore,
\[
\begin{aligned}
\E T&=1+\int_0^1f(x)\dd x\\
&=A\sin1+\cos1
=\boxed{\sec1+\tan1}.
\end{aligned}
\]
The trigonometric arguments are in radians.
